% ============================================================
% Cover Letter — Environmental Modelling & Software
% Manuscript: Ceará Digital Twin for Wildfires (EN)
% ============================================================
\documentclass[a4paper,12pt]{letter}

\usepackage[utf8]{inputenc}
\usepackage[T1]{fontenc}
\usepackage[english]{babel}
\usepackage{hyperref}
\usepackage{geometry}
\geometry{margin=2.5cm}

\signature{Nauber Gois\\On behalf of all co-authors}
\address{%
  Universidade de Fortaleza (UNIFOR)\\
  Fortaleza, Ceará, Brazil\\
  \texttt{naubergois@gmail.com}
}

\begin{document}

\begin{letter}{%
  Editorial Office\\
  \textit{Environmental Modelling \& Software}\\
  Elsevier
}

\opening{Dear Editor,}

We submit the manuscript entitled:

\begin{quote}
\textbf{Ceará Digital Twin for Wildfires: An Open-Source Agentic AI Platform\\
Integrating Near Real-Time Satellite Data and LangGraph Orchestration}
\end{quote}

for consideration as a research article in \textit{Environmental Modelling \& Software}.

\medskip
\textbf{Summary.}\ 
Wildfire monitoring in Brazil still depends largely on centralized systems with multi-hour latency between satellite overpass and public data availability. We present an open-source, agentic wildfire monitoring platform for Ceará that integrates NASA FIRMS, optional GOES-16, INPE BDQueimadas, and Open-Meteo within a ten-node LangGraph production pipeline (ingestion, validation, parallel geospatial/climate/GOES analysis, evidence fusion, risk classification, ReAct diagnosis, and three-class alerts).

The manuscript combines (i)~operational software engineering for environmental modelling workflows, (ii)~a three-class alert framework (NO/UNCERTAIN/YES) with explicit precision--recall trade-offs, and (iii)~NeKo-PIGNN (Neural Koopman Operator + Physics-Informed GNN) as an offline research module for physics-regularized prediction benchmarks.

\medskip
\textbf{Key Results.}
\begin{itemize}
  \item \textbf{Operational alerts:} 82--92\% YES precision on a 20-day holdout; bootstrap 95\% CI [68\%, 96\%] for XGBoost YES precision.
  \item \textbf{Extended INPE validation:} 84.2\% precision and 91.8\% recall over the 2025 dry season (184 days, 24{,}573 hotspots).
  \item \textbf{NeKo-PIGNN:} $R^2=0.972$ on synthetic benchmarks; real-data results are reported separately from operational alert metrics to avoid conflation.
  \item \textbf{Reproducibility:} MIT-licensed code, experiment scripts (Exp-B v1--v9 and EXP-ROBUST-001--004), and pinned result tables on GitHub (\url{https://github.com/naubergois/ceara-queimadas}).
\end{itemize}

\medskip
\textbf{Fit with the Journal.}\
This work addresses software for environmental modelling and decision support: modular agent orchestration, reproducible experiment pipelines, integration of heterogeneous geospatial/climate data streams, and transparent reporting of limitations (GOES-16 not yet in continuous production; seasonal precision variation). We believe it will interest readers working on wildfire monitoring systems, operational AI for civil defense, and open environmental software.

\medskip
\textbf{Suggested Reviewers.}
\begin{enumerate}
  \item Prof.\ Steven L.\ Brunton --- University of Washington (Koopman operators, data-driven dynamical systems)
  \item Prof.\ George Em Karniadakis --- Brown University (physics-informed machine learning)
  \item Prof.\ Alan A.\ Ager --- USDA Forest Service (wildfire spread modelling and operational tools)
\end{enumerate}

\medskip
\textbf{Declarations.}\
All five authors (Nauber Gois, Alexandre Lobo, Vladia Celia Monteiro Pinheiro, Daniel Valente de Macedo, Adriano Bessa Albuquerque) approved the manuscript. We declare no competing interests. The study uses publicly available satellite and meteorological data; no human subjects were involved.

\closing{Sincerely,}

\end{letter}
\end{document}
